\documentclass[12pt,a4paper]{article}

\usepackage[spanish]{babel}
\usepackage[utf8]{inputenc}
\usepackage[T1]{fontenc}
\usepackage{geometry}
\usepackage{enumitem}
\usepackage{hyperref}
\geometry{margin=2.5cm}

\title{Propuesta Funcional de la App \\ Gestión de Pagos y Validación de Comprobantes \\ para Viajes Grupales}
\author{}
\date{\today}

\begin{document}

\maketitle

\section{Alcance del sistema}
La app está orientada a la gestión de pagos y validación de comprobantes de viajes grupales administrados por una agencia.
Los administradores son el dueño de la agencia y/o empleados autorizados.
Los usuarios son padres/tutores que pagan el viaje de sus hijos.

\section{¿Qué es la app?}
Es una plataforma web para administrar pagos de viajes grupales.
La app no cobra dinero directamente: funciona como un registro digital donde los padres/tutores informan pagos y suben comprobantes, y la administración de la agencia valida esos comprobantes.

\subsection*{Incluye}
\begin{itemize}[leftmargin=*]
    \item Panel de usuarios (padres/tutores) con vista de todos sus viajes, cuotas y estados de pago.
    \item Panel de administración (dueño/empleados de la agencia).
    \item \textbf{Plantilla interna tipo Excel integrada en la app} (sin depender de Excel externo para operar).
    \item Descarga de plantilla en archivo Excel (\texttt{.xlsx}).
    \item Gestión segura de comprobantes (JPG/PDF).
\end{itemize}

\section{Aclaración clave sobre la plantilla interna}
La plantilla interna es una implementación propia de la app con comportamiento tipo Excel (grilla tabular editable, filtros y ordenamiento), integrada directamente en el sistema.
\textbf{No es un archivo Excel externo embebido ni una dependencia de Google Sheets/Microsoft Excel para trabajar día a día.}
La operación diaria se realiza dentro de la app y, adicionalmente, se puede exportar la información a \texttt{.xlsx} cuando se necesite compartir o archivar.

\section{¿Qué resuelve?}
\begin{itemize}[leftmargin=*]
    \item Centraliza la información de cuotas de cada viaje.
    \item Permite saber qué está pago, pendiente o vencido.
    \item Ordena la validación de comprobantes.
    \item Facilita el seguimiento administrativo con una plantilla clara y exportable.
    \item Identifica cuotas próximas a vencer y cuotas vencidas para facilitar el seguimiento.
\end{itemize}

\section{Roles del sistema}

\subsection{Administrador (Agencia)}
\begin{itemize}[leftmargin=*]
    \item Crear, editar y gestionar viajes.
    \item Definir por viaje:
    \begin{itemize}
        \item monto total,
        \item cantidad de cuotas,
        \item día de vencimiento,
        \item días previos para estado ``amarillo'',
        \item \textbf{aplicar retroactividad a nuevos inscriptos (sí/no)}.
    \end{itemize}
    \item Asignar/desasignar familias (usuarios) a viajes (individual y masivo opcional).
    \item Ver cuentas de cada usuario.
    \item Revisar comprobantes y aceptar/rechazar con observación.
    \item Usar la plantilla interna tipo Excel del viaje.
    \item Descargar la plantilla en Excel (\texttt{.xlsx}).
    \item Ver y descargar comprobantes con permisos de administrador.
\end{itemize}

\subsection{Usuario (Padre/Tutor)}
\begin{itemize}[leftmargin=*]
    \item Al iniciar sesión, ver inmediatamente el estado de pago de \textbf{todos sus viajes activos}.
    \item Por cada viaje, visualizar \textbf{cada cuota individualmente} con su estado representado por color:
    \begin{itemize}
        \item \textcolor{green}{\textbf{Verde}}: cuota pagada y validada por administración, o con mucho tiempo hasta el vencimiento.
        \item \textcolor{orange}{\textbf{Amarillo}}: cuota próxima al vencimiento, o con comprobante pendiente de aprobación por el administrador.
        \item \textcolor{red}{\textbf{Rojo}}: cuota vencida sin regularizar, o comprobante rechazado por el administrador.
    \end{itemize}
    \item Ver su estado de cuenta:
    \begin{itemize}
        \item deuda actual,
        \item cuotas pagadas y pendientes,
        \item próximos vencimientos,
        \item \textbf{total actualizado a pagar} por cuota.
    \end{itemize}
    \item Cargar comprobantes de pago (JPG/PDF) informando:
    \begin{itemize}
        \item importe,
        \item fecha de pago,
        \item método de pago.
    \end{itemize}
    \item Ver el estado de sus comprobantes (pendiente de revisión, aprobado o rechazado).
    \item Si un comprobante es rechazado, ver la observación del administrador y poder enviar uno nuevo.
\end{itemize}

\section{Semáforo de pagos}
Cada cuota se representa visualmente con un color. Los estados internos del sistema se traducen a colores de la siguiente manera:

\begin{center}
\begin{tabular}{|l|l|l|}
\hline
\textbf{Color mostrado} & \textbf{Estado interno} & \textbf{Significado} \\
\hline
\textcolor{green}{\textbf{Verde}} & GREEN & Cuota pagada y validada por admin \\
\hline
\textcolor{orange}{\textbf{Amarillo}} & YELLOW & Próxima al vencimiento o pendiente de revisión \\
\hline
\textcolor{red}{\textbf{Rojo}} & RED & Cuota vencida sin pagar \\
\hline
\textcolor{red}{\textbf{Rojo}} & RETROACTIVE (rechazado) & Comprobante rechazado por admin \\
\hline
\end{tabular}
\end{center}

\textbf{Regla de prioridad de color:} si una cuota tiene un comprobante rechazado, se muestra en rojo independientemente de su fecha de vencimiento.
Si tiene un comprobante pendiente de aprobación, se muestra en amarillo.
El texto que ve el usuario nunca expone los nombres técnicos internos (GREEN, YELLOW, RED, RETROACTIVE): siempre se muestran etiquetas legibles en español.

\section{Total exigible de una cuota}
El total exigible de una cuota normal deriva de su capital. Una cuota vencida puede mostrarse en rojo por estar vencida e impaga, pero su estado no agrega un importe automático al total.

\section{Notificaciones automáticas}
Cuando una cuota pasa a estado amarillo, la app envía un email recordatorio al usuario con:
\begin{itemize}[leftmargin=*]
    \item viaje,
    \item cuota/período,
    \item monto,
    \item fecha límite de pago,
    \item recordatorio del vencimiento.
\end{itemize}

\section{Plantilla interna tipo Excel (una sola por viaje)}
Cada viaje tiene una única plantilla administrativa integrada en la app.
La plantilla consolida todos los usuarios del viaje en una sola vista y agrega tantas columnas de cuota como corresponda al plan de pagos definido.

\subsection*{Estructura}
\begin{itemize}[leftmargin=*]
    \item Columnas de datos del usuario: Nombre, Apellido, Teléfono, Email, Alumno, Colegio, Curso.
    \item Columnas dinámicas de cuotas: \texttt{Cuota 1, Cuota 2, ..., Cuota N}.
    \item Cada columna de cuota refleja:
    \begin{itemize}
        \item estado representado por color (verde/amarillo/rojo),
        \item capital de cuota,
        \item retroactivo (si aplica),
        \item total exigible.
    \end{itemize}
\end{itemize}

\subsection*{Funciones de la plantilla}
\begin{itemize}[leftmargin=*]
    \item filtros,
    \item ordenamiento,
    \item búsqueda,
    \item navegación por páginas.
\end{itemize}

\section{Regla de retroactividad (Opcional)}
Al momento de crear un viaje, la administración puede definir si el cobro retroactivo aplica o no para ese grupo.

\subsection*{Criterio operativo}
\begin{itemize}[leftmargin=*]
    \item \textbf{Si está inactiva:} El usuario que ingresa tarde solo es responsable de las cuotas con fecha de vencimiento posterior a su alta.
    \item \textbf{Si está activa:} Las cuotas anteriores al momento del alta se consideran deuda retroactiva y se suman al importe de la primera cuota pendiente.
\end{itemize}

\section{Estructura de cuotas por usuario}
Para cada viaje, el sistema genera una cuota por cada período para cada usuario asignado.
Cada usuario tiene su propio conjunto de cuotas dentro de cada viaje.

\section{Exportación a Excel (\texttt{.xlsx})}
Desde administración se puede descargar la plantilla del viaje en formato \texttt{.xlsx}, con los datos actualizados y formato de lectura claro para compartir o archivar.

\section{Seguridad de archivos}
\begin{itemize}[leftmargin=*]
    \item Los comprobantes no quedan públicos en internet.
    \item Solo usuarios autorizados (administradores de la agencia) pueden acceder.
    \item El acceso a archivos se controla desde la app con autenticación.
\end{itemize}

\section{Reglas operativas importantes}
\begin{itemize}[leftmargin=*]
    \item No se puede marcar un pago como aceptado sin comprobante asociado.
    \item Si se rechaza un comprobante, se registra motivo/observación y la cuota se muestra en rojo.
    \item El estado del semáforo se actualiza automáticamente según fecha y validación.
    \item Cada usuario ve solo sus datos; administración ve el conjunto completo.
    \item Se puede agregar usuarios al viaje en cualquier momento, aplicando la regla de retroactividad según la configuración del viaje.
    \item \textbf{Los colores mostrados al usuario nunca exponen nombres técnicos internos}: siempre se usan etiquetas en español claras y comprensibles.
\end{itemize}

\section{Resultado esperado}
Una app completa para la agencia que permita operar viajes con control de cuotas por usuario, validación de comprobantes, alertas automáticas, seguimiento en plantilla interna tipo Excel integrada y exportación a Excel, con reglas claras de vencimiento y trazabilidad administrativa.

\end{document}
